%% kapa: exam=1 round=26 year=2015 period=2 form=A subject=appraisal_law subject_ko=부동산관계법규 seq=19 number=19
%% source: Q-Net 감정평가사 기출문제 (cst00309) artlSeq=5138001 file=extracted/1008/5138001/제26회_감정평가사 1차_기출/제26회_감정평가사 1차_기출/2교시_감평_A형.pdf
%% conversion: kordoc 4.24.51 -> tools/gen1.py (auto); review=proofread-sonnet; flags=none
\begin{question}{19}
부동산가격공시 및 감정평가에 관한 법령상 지가 및 주택가격의 공시에 관한 설명으로 옳지 않은 것은?
\begin{choices}
  \item 표준주택가격의 공시에는「국토의 계획 및 이용에 관한 법률」등에 의한 토지의 용도제한이 포함되어야 한다.
  \item 개별주택가격의 공시에는 표준주택의 대지면적 및 형상이 포함되어야 한다.
  \item 표준지공시지가에 대하여 이의가 있는 자는 표준지공시지가의 공시일부터 30일 이내에 서면으로 국토교통부장관에게 이의를 신청할 수 있다.
  \item 시장ㆍ군수 또는 구청장은 국토교통부장관이 정한 개별공시지가의 조사ㆍ산정지침에 따라 개별공시지가를 조사ㆍ산정하여야 한다.
  \item 국토교통부장관이 따로 정하지 아니한 경우 표준주택가격의 공시기준일은 1월 1일로 한다.
\end{choices}
\end{question}
\answer{2}
